\section{\scalebox{0.84}[1]{Proposed MBA-RL for Charging Balancing Control}}
\newrev{As shown in Fig.~\ref{fig:MARL_Schematic}, each cell has a local state and an independently regulated charging current, while all cells share the objective of fast charging and pack-level SOC balancing. The resulting cooperative task uses a discrete current set and requires decentralized decisions during execution. QMIX is therefore adopted because its centralized mixing network learns a joint action value during training, while the monotonic value decomposition preserves decentralized greedy action selection \cite{26}. The QMIX controller is integrated with multiple SPMe-based cell models to form the MBA-RL framework described below.}


\subsection{MBA-RL Charging Control Framework}
\newrev{The schematic of MBA-RL is presented in Fig.~\ref{fig:schematic_method}. MBA-RL comprises shared agent networks, a mixing network, and hypernetworks. For the $i$-th agent, a recurrent Q-network combines the current observation $o_t^i$ with the previous executed action $\tilde{a}_{t-1}^i$ to form the action-observation history $\tau_t^i$ and generate the local action value $Q_i$. The mixing network combines $Q_1,Q_2,\ldots,Q_n$ into the joint action value $Q_{\mathrm{tot}}$, and the hypernetworks generate its state-dependent weights. Target agent and mixing networks provide temporal-difference targets during training. The shared agent network receives the local voltage, temperature, and SOC together with the previous-action encoding and an agent identifier.}

\newrev{We formulate the charging balancing problem as a Dec-POMDP \cite{27} described by the tuple}
{\color{red}
\begin{equation*}
\begin{aligned}
G=\langle &\mathcal{N},\mathcal{S},\{\mathcal{A}_i\}_{i=1}^{n},\mathcal{P},\mathcal{R},\\[-2pt]
&\{\mathcal{O}_i\}_{i=1}^{n},\mathcal{O},\gamma\rangle .
\end{aligned}
\end{equation*}
}
\newrev{Here, $\mathcal{N}=\{1,\ldots,n\}$ is the agent set, $\mathcal{S}$ is the global state space, $\mathcal{A}_i$ and $\mathcal{O}_i$ are the action and observation spaces of agent $i$, $\mathcal{P}$ is the transition kernel, $\mathcal{R}$ is the shared reward function, $\mathcal{O}$ is the observation function, and $\gamma$ is the discount factor.}

\subsubsection{\newrev{State and Observation Spaces}} \newrev{The global state $\mathbf{s}_t\in\mathcal{S}$ is formed by the observations of all agents at time step $t$:}
\begin{equation}
   \mathbf{s}_t = \left[ o_t^1, o_t^2, \dots, o_t^n \right]^{\mathrm{T}}
\end{equation}
\newrev{For agent $i$, the local observation $o_t^i\in\mathcal{O}_i$ consists of the cell voltage $V_i(t)$, average temperature $T_i(t)$, and $SOC_i(t)$. The observation function $\mathcal{O}$ maps the battery-pack state and joint action to the local observations available to the agents.}

\subsubsection{Action Space} \newrev{Each agent selects a requested current from the discrete candidate set $\mathcal{A}_{\mathrm{cand}}=\{0,0.5,\ldots,7.5\}\,\mathrm{A}$. Before action selection, a state-dependent availability mask removes commands that conflict with the current sampled state. Only zero current is available when $SOC_i\geq0.95$, while currents above 1 A are excluded when $V_i\geq V_{\max}$ or $T_i\geq T_{\max}$. The mask restricts the actions considered by the policy but does not predict constraint satisfaction over the next interval; therefore, the selected action is subsequently processed by the safety-constrained execution mechanism.}
\begingroup
\setlength{\abovedisplayskip}{3pt}
\setlength{\belowdisplayskip}{3pt}
\begin{equation}
\newrev{\resizebox{0.91\columnwidth}{!}{$\displaystyle
\mathbf{a}_t=\left[a_t^1,a_t^2,\dots,a_t^n\right]^{\mathrm{T}},\quad
\mathcal{A}(\mathbf{s}_t)=\prod_{i=1}^{n}\mathcal{A}_i(\mathbf{s}_t),\quad
\mathcal{A}_i(\mathbf{s}_t)\subseteq\mathcal{A}_{\mathrm{cand}}.$}}
\end{equation}
\endgroup
\subsubsection{Reward Function} \newrev{The reward combines charging progress, pack-level SOC balance, and voltage and temperature penalties:}
\begin{equation}
    r_t = r_{\text{time}} + r_{\text{bal}} + r_{\text{safety}}
\end{equation}
\newrev{The time penalty is $r_{\text{time}}=-\lambda_{\text{time}}$ before the terminal condition is reached.}

Then $r_{\text{bal}}$ can be calculated as follows:
\begin{equation}
  r_{\text{bal}} = 
\begin{cases}
 -\lambda_{\text{bal}}(\sigma - \beta), & \text{if } \sigma \geq \beta \\
0, & \text{otherwise}
\end{cases}  
\end{equation}
where $\sigma$ denotes the standard deviation of the SOC across all cells in the battery pack, $\beta$ indicates the predefined threshold, and $r_{\text{bal}}$ characterizes the imbalance effects in the battery pack. 

As shown in \eqref{safty condition}, $r_{\text{safety}}$ is divided into voltage safety constraints and temperature safety constraints:
\begin{equation}
    r_{\text{safety}} = r_{\text{volt}} + r_{\text{temp}}
\end{equation}
\begin{equation}
    \quad r_{\text{volt}} = \sum_{i=1}^{n} \delta_{v,i}, \quad r_{\text{temp}} = \sum_{i=1}^{n} \delta_{t,i}
\end{equation}
where $\delta_{v,i}$ and $\delta_{t,i}$ represent the overvoltage and overtemperature penalties of the $i$-th cell, respectively:
\begin{equation}
    \delta_{v,i} = 
\begin{cases}
 -\lambda_v(V_i(t) - V_{\text{max}}), & \text{if } V_i(t) \geq V_{\text{max}} \\
0, & \text{otherwise}
\end{cases}
\end{equation}

\begin{equation}
  \delta_{t,i} = 
\begin{cases}
 -\lambda_T(T_i(t) - T_{\text{max}}), & \text{if } T_i(t) \geq T_{\text{max}} \\
0, & \text{otherwise}
\end{cases}  .
\end{equation}
\subsubsection{\newrev{Transition and Termination}} \newrev{The safety screen is treated as part of the action-execution interface. Given the requested joint action $\mathbf{a}_t$, it produces the executed action $\tilde{\mathbf{a}}_t$, and the transition kernel $\mathcal{P}(\mathbf{s}_{t+1}\mid\mathbf{s}_t,\mathbf{a}_t)$ is induced by this screening operation and the battery dynamics. The variable $done$ is defined separately as the episode-termination indicator.}

\begin{algorithm}[t]
    \caption{MBA-RL charging control framework}
    \label{alg:QMIX}
    
    \begin{algorithmic}[1]
        \STATE Initialize $\theta, \theta', \bm{\theta}, \bm{\theta}'$ randomly
        \STATE Set the learning rate $\alpha$, number of episodes $N_{episode}$, and replay buffer $\mathcal{D}$
        \STATE $\theta' \leftarrow \theta, \bm{\theta}' \leftarrow \bm{\theta}$
        \FOR{$episode =$ 1 to $N_{episode}$}
            \STATE $t\leftarrow 0$, $\tau_t^i \leftarrow \emptyset$, and observe the initial state $\mathbf{s}_t$

            // \textbf{Data Collection}
            \WHILE{$\neg done$}
                \FOR{the $i$-th agent}
                    \STATE \newrev{$\tau_t^i \leftarrow \tau_{t-1}^i \cup \{(o_t^i, \tilde{a}_{t-1}^i)\}$}
                    \STATE \newrev{Update $\epsilon$ according to the exploration schedule}
                    \STATE Select the requested action $a_t^i$ according to \eqref{epsilon-greedy}
                \ENDFOR
                \STATE \newrev{Screen $\mathbf{a}_t$ to obtain $\tilde{\mathbf{a}}_t$, execute $\tilde{\mathbf{a}}_t$, observe $r_t$, and get $\mathbf{s}_{t+1}$}
                \STATE \newrev{$\mathcal{D} \leftarrow \mathcal{D} \cup \{ \langle\mathbf{s}_t, \mathbf{m}_t, \mathbf{a}_t, \tilde{\mathbf{a}}_t, r_t, \mathbf{s}_{t+1}, \mathbf{m}_{t+1}, done\rangle \}$}
            \ENDWHILE
            
            // \textbf{Network Update}
            \IF{$|\mathcal{D}| > N$}
                \STATE $B \leftarrow$ Randomly sample data from $N$ episodes in $\mathcal{D}$
                \FOR{each time step $t$ in each episode in batch $B$}
                    \STATE Calculate $Q_{tot}$ by \eqref{Q-value} based on the state $\mathbf{s}_t$
                    \STATE Calculate $y^{\text{tot}}$ by \eqref{target Q} based on the state $\mathbf{s}_{t+1}$
                \ENDFOR
                \STATE Under the conditions given in \eqref{QMIX_constraints} and \eqref{QMIX_constraints_1}, utilize the RMSprop  optimizer \cite{zou2019sufficient} to update $\theta$ and $\bm{\theta}$ by minimizing the loss function: $\left( y_{\text{tot}} - Q_{\text{tot}} \right)^2$
            \ENDIF
            \IF{$episode$ \% $N_t$ == 0}
                \STATE $\theta' \leftarrow \theta$, $\bm{\theta}' \leftarrow \bm{\theta}$ 
            \ENDIF
        \ENDFOR
    \end{algorithmic}
\end{algorithm}

\subsection{Safety-Constrained Execution}
\newrev{At each decision step, MBA-RL generates a requested current for every cell. Starting from this request, the safety screen evaluates successively lower candidate currents with the nominal SPMe over the next decision interval and a subsequent zero-current interval. Prediction corrections derived from the observed-model mismatch are added to the predicted SOC, voltage, and temperature peaks. The largest candidate satisfying the margins specified in Table~\ref{tab:experimental_settings} is executed. If no feasible candidate exists for any cell, the episode terminates without advancing the battery model. The same screen is used during training and execution. It provides a sampled finite-horizon check rather than a continuous-time safety guarantee.}

\subsection{Balancing Mechanism of MBA-RL}

\begin{figure}[t]
    \centering
    \includegraphics[width=0.95\linewidth]{experimental_platform.png}
    \caption{Schematic of the experimental platform.}
    \label{fig:experiment_Schematic}
\end{figure}

\newrev{MBA-RL uses global information during centralized training so that each agent accounts for the charging-time objective in \eqref{eqn:balancing objective} together with the pack-level balancing requirement in \eqref{terminal balance}.}

In MBA-RL, the $i$-th agent aims to optimize its local Q-value, denoted as $Q_i(\tau_t^i, a_t^i|\theta)$, to maximize the global Q-value $Q_{\text{tot}}$. This involves balancing the agent's own goals of safety and fast charging while simultaneously contributing to the rapid achievement of balance within the battery pack. Each agent chooses its current action $a_t^i$ using the $\epsilon$-greedy strategy, which can be described as follows:
\begin{equation}\label{epsilon-greedy}
a_t^i =
\begin{cases}
\text{random available action} & \text{with probability } \epsilon \\
\arg\max_{a_t^i} Q_i(\tau_t^i, a_t^i|\theta) & \text{with probability } 1 - \epsilon
\end{cases}
\end{equation}
where $Q_i(\tau_t^i, a_t^i|\theta)$ is computed based on $\tau_t^i$ and $a_t^i$. \rev{All agents share a common set of parameters $\theta$. A one-hot agent identifier is appended to each network input so that the shared network can represent cell-specific action values.}

\newrev{The mixing network takes the outputs of all agent networks as inputs and combines them monotonically to generate $Q_{tot}(\boldsymbol{\tau}_t, \mathbf{a}_t|\bm{\theta})$:}
\begin{equation}\label{Q-value}
    Q_{tot}(\boldsymbol{\tau}_t, \mathbf{a}_t|\bm{\theta}) = f(Q_1(\tau_t^1, a_t^1|\theta), ..., Q_n(\tau_t^n, a_t^n|\theta))
\end{equation}
where $\boldsymbol{\tau}_t$ denotes the joint action-observation history at the time step $t$, and $\bm{\theta}$ represents the parameters of the mixing network. 

The final objective of MBA-RL is to maximize the global Q-value $Q_{tot}(\boldsymbol{\tau}_t, \mathbf{a}_t|\bm{\theta})$. \newrev{The mixing network enforces a monotonicity constraint between the joint and local action values, so decentralized greedy action selection is consistent with joint greedy selection:}
\begin{equation}\label{QMIX_constraints}
\arg\max_{\mathbf{a}_t} Q_{tot}(\boldsymbol{\tau}_t, \mathbf{a}_t|\bm{\theta}) =
\begin{pmatrix}
\arg\max_{a_t^1} Q_1(\tau_t^1, a_t^1|\theta) \\
\vdots \\
\arg\max_{a_t^n} Q_n(\tau_t^n, a_t^n|\theta)
\end{pmatrix}.
\end{equation}
Eq. \eqref{QMIX_constraints} enables each agent to greedily select the action that maximizes $Q_i(\tau_t^i, a_t^i|\theta)$ for decentralized execution. 

\newrev{The hypernetworks receive the state $s_t$ and generate nonnegative weights for the mixing network, enforcing}
\begin{equation}\label{QMIX_constraints_1}
\frac{\partial Q_{tot}}{\partial Q_i} \geq 0, \, i = 1, 2, ..., n.
\end{equation}

\newrev{The target agent and mixing networks are initialized from their corresponding online networks and synchronized every $N_t$ training episodes. Considering the terminal flag and the availability mask at the next state, the global target Q-value is computed as}
\begin{equation}\label{target Q}
\newrev{y_t^{\text{tot}} = r_t + \gamma(1-d_t) \max_{\mathbf{a}_{t+1}\in\mathcal{A}(\mathbf{s}_{t+1})} Q_{\text{tot}}(\boldsymbol{\tau}_{t+1}, \mathbf{a}_{t+1}|\bm{\theta}'),}
\end{equation}
\newrev{where $d_t$ denotes the terminal flag.}

\newrev{Complete episodes are stored in the replay buffer $\mathcal{D}$ together with the observations, availability masks, requested and executed actions, rewards, and terminal flags. During training, a batch of episodes is randomly sampled for network updates. The Q-function is indexed by the requested action selected by the policy, while the executed action is included in the subsequent action-observation history to account for intervention by the safety screen.}

In summary, the details of the MBA-RL charging control framework are described in Algorithm 1, where $|\cdot|$ denotes the cardinality of a set.

\begin{figure}[t]
    \centering

    \begin{subfigure}[t]{0.24\textwidth} 
        \centering
        \includegraphics[width=\textwidth]{aging_voltage.png}
        \caption{}
    \end{subfigure}
    \hfill
    \begin{subfigure}[t]{0.24\textwidth}
        \centering
        \includegraphics[width=\textwidth]{aging_temperature.png} 
        \caption{}
    \end{subfigure}

    \caption{Voltage and temperature curves of the three cells at 1 C discharging current. (a) Voltage curve. (b) Temperature curve.}
    \label{fig:three_aging}
\end{figure}

\begin{figure}[t]
    \centering

    \begin{subfigure}[t]{0.24\textwidth} 
        \centering
        \includegraphics[width=\textwidth]{cell1_voltage_fit.png}
        \caption{}
    \end{subfigure}
    \hfill
    \begin{subfigure}[t]{0.24\textwidth}
        \centering
        \includegraphics[width=\textwidth]{cell1_temperature_fit.png} 
        \caption{}
    \end{subfigure}

    \caption{Parameter identification results for Cell 1. (a) Electrochemical model parameter identification. (b) Thermal model parameter identification.}
    \label{battery1}
\end{figure}

\begin{figure}[t]
    \centering

    \begin{subfigure}[t]{0.24\textwidth} 
        \centering
        \includegraphics[width=\textwidth]{cell2_voltage_fit.png}
        \caption{}
    \end{subfigure}
    \hfill
    \begin{subfigure}[t]{0.24\textwidth}
        \centering
        \includegraphics[width=\textwidth]{cell2_temperature_fit.png} 
        \caption{}
    \end{subfigure}

    \caption{Parameter identification results for Cell 2. (a) Electrochemical model parameter identification. (b) Thermal model parameter identification.}
    \label{battery2}
\end{figure}

\begin{figure}[t]
    \centering

    \begin{subfigure}[t]{0.24\textwidth} 
        \centering
        \includegraphics[width=\textwidth]{cell3_voltage_fit.png}
        \caption{}
    \end{subfigure}
    \hfill
    \begin{subfigure}[t]{0.24\textwidth}
        \centering
        \includegraphics[width=\textwidth]{cell3_temperature_fit.png} 
        \caption{}
    \end{subfigure}

    \caption{Parameter identification results for Cell 3. (a) Electrochemical model parameter identification. (b) Thermal model parameter identification.}
    \label{battery3}
\end{figure}
